
\documentclass[11pt]{article}
\usepackage[margin=1in]{geometry}

% Core packages
\usepackage{amsmath,amssymb}
\usepackage{graphicx}
\usepackage{xcolor}
\usepackage{needspace}
\usepackage{tikz-cd}
\usepackage{multicol}
\newcommand{\termnote}[2]{\underbrace{#1}_{\begingroup\renewcommand{\arraystretch}{0.65}\begin{array}{c}#2\end{array}\endgroup}}
\newenvironment{assignmentinstruction}{\par\begingroup\color{blue}}{\par\endgroup}

% Paragraphs
\setlength{\parindent}{0pt}
\setlength{\parskip}{1\baselineskip}

\title{Analytical Dynamics Project 2: Internal Resonance in a Nanomechanical Oscillator}
\author{Author Name}
\date{\today}

\begin{document}
\maketitle

\section*{Introduction}

This project studies a reduced-order model of a hybrid, multi-span nanomechanical resonator operating near primary resonance. The physical system consists of inner and outer beam fields connected at compliant hinges. Magnetic actuation is applied near the midpoint, while electrostatic forces act through capacitors near the outer edges. A Galerkin assumed-mode reduction represents the coupled beam dynamics using two generalized coordinates, $y_1$ and $y_2$; these are reduced model coordinates rather than direct pointwise beam deflections. Figure~\ref{fig:nanoresonator-setup} illustrates the physical setup and its reduced two-field representation. Motivated by the model in Alcorta et al.~\cite{alcorta2021}, we analyze equilibrium stability, natural frequencies, internal resonance, and the forced response using the method of multiple scales.

\begin{figure}[!t]
\centering
\includegraphics[page=3,trim=105pt 235pt 105pt 330pt,clip,width=0.9\textwidth]{\detokenize{Assignment/Original problem/clapper_27dec21_agbg.pdf}}
\caption{Physical nanomechanical resonator and reduced two-field scheme. Reproduced from Fig.~1 of Alcorta et al.~\cite{alcorta2021}.}
\label{fig:nanoresonator-setup}
\end{figure}

\section{Problem Formulation}

The \underline{nondimensional equations of motion} are:
{\small
\begin{equation}
\begin{aligned}
\termnote{y_{1,\tau\tau}}{\text{\scriptsize inertia}}
+ \termnote{\zeta y_{1,\tau}}{\text{\scriptsize linear}\\\text{\scriptsize damping}}
+ \termnote{\omega_1^2 y_1}{\text{\scriptsize linear}\\\text{\scriptsize stiffness}}
+ \termnote{\alpha_{11} y_1^3}{\text{\scriptsize cubic}\\\text{\scriptsize stiffness}}
+ \termnote{\alpha_{12} y_1 y_2^2}{\text{\scriptsize cross}\\\text{\scriptsize coupling}}
+ \termnote{\eta_{11} y_1^2 y_{1,\tau}}{\text{\scriptsize nonlinear}\\\text{\scriptsize damping}}
+ \termnote{\eta_{12} y_1 y_2 y_{2,\tau}}{\text{\scriptsize cross}\\\text{\scriptsize damping}}
&= \termnote{f_1}{\text{\scriptsize magnetic}\\\text{\scriptsize forcing}}
\mkern-8mu+\mkern-8mu \termnote{g_1}{\text{\scriptsize electrostatic}\\\text{\scriptsize force}},\\
\termnote{y_{2,\tau\tau}}{\text{\scriptsize inertia}}
+ \termnote{\zeta y_{2,\tau}}{\text{\scriptsize linear}\\\text{\scriptsize damping}}
+ \termnote{\omega_2^2 y_2}{\text{\scriptsize linear}\\\text{\scriptsize stiffness}}
+ \termnote{\alpha_{21} y_1^2 y_2}{\text{\scriptsize cross}\\\text{\scriptsize coupling}}
+ \termnote{\alpha_{22} y_2^3}{\text{\scriptsize cubic}\\\text{\scriptsize stiffness}}
+ \termnote{\eta_{11} y_1^2 y_{1,\tau}}{\text{\scriptsize nonlinear}\\\text{\scriptsize damping}}
+ \termnote{\eta_{12} y_1 y_2 y_{2,\tau}}{\text{\scriptsize cross}\\\text{\scriptsize damping}}
&= \termnote{g_2}{\text{\scriptsize electrostatic}\\\text{\scriptsize force}}.
\end{aligned}
\end{equation}
}
Here, $y_1$ and $y_2$ are the generalized coordinates of the reduced model, and $\tau$ is nondimensional time. The parameter $\zeta$ represents linear damping; $\omega_1$ and $\omega_2$ are the linear frequency parameters; the $\alpha_{ij}$ describe cubic stiffness and intermodal coupling; and the $\eta_{ij}$ describe nonlinear damping. The magnetic forcing has amplitude $\mu$ and frequency $\Omega$, while $\beta$ sets the electrostatic force scale. The coordinates and parameters are nondimensional, and all parameters are positive as specified in the project appendix.

\par\medskip
The \underline{magnetic force} is:
{\small
\begin{equation}
f_1 = \termnote{\mu\sin(\Omega\tau)}{\text{\scriptsize magnetic}\\\text{\scriptsize input}}
\end{equation}
}

\par\medskip
The \underline{electrostatic forces} are:
{\small
\begin{equation}
g_{1,2} = \beta\bigl(\termnote{\frac{1}{[1-(y_1+y_2)]^2}}{\text{\scriptsize gap}\\\text{\scriptsize one}}
\pm \termnote{\frac{1}{[1-(y_1-y_2)]^2}}{\text{\scriptsize gap}\\\text{\scriptsize two}}\bigr)
\end{equation}
}
The upper sign in the electrostatic-force expression defines $g_1$, and the lower sign defines $g_2$. The normalized gaps must remain positive, so the model is physically defined only when $1-(y_1+y_2)>0$ and $1-(y_1-y_2)>0$.

\Needspace{18\baselineskip}
In the two-mode Galerkin description, the \underline{displacement-field expansion} is:
\begin{equation}
\begin{aligned}
u^A(s,\tau)&=\phi_1^A(s)q_1(\tau)+\phi_2^A(s)q_2(\tau),\\
u^B(s,\tau)&=\phi_1^B(s)q_1(\tau)+\phi_2^B(s)q_2(\tau),\\
y_1&=\frac{\phi_1^B(1)q_1}{\gamma},\qquad
y_2=\frac{\phi_2^B(1)q_2}{\gamma},\\
\frac{u^B(1,\tau)}{\gamma}&=y_1+y_2,\qquad
\frac{u^B(-1,\tau)}{\gamma}=y_1-y_2.
\end{aligned}
\end{equation}
Here, $s$ is the nondimensional position along the beam, $\phi_i^{A,B}(s)$ are spatial mode shapes, and $q_i(\tau)$ are their time-varying modal amplitudes. The first mode is symmetric and the second is antisymmetric. Thus, $y_1$ and $y_2$ are scaled modal amplitudes—not mode shape functions themselves—and the displacement at any point is reconstructed by the first two lines. The last line gives the normalized outer-edge displacements and explains the combinations in the electrostatic-force denominators.

\begin{figure}[!t]
\centering
\includegraphics[page=6,trim=105pt 540pt 105pt 110pt,clip,width=0.9\textwidth]{\detokenize{Assignment/Original problem/clapper_27dec21_agbg.pdf}}
\caption{First symmetric mode (left) and second antisymmetric mode (right); the two curves in each panel show the inner and outer beam fields. Reproduced from Fig.~2 of Alcorta et al.~\cite{alcorta2021}.}
\label{fig:mode-shapes}
\end{figure}

\section{Equilibrium Analysis}
\subsection{Equilibrium and Stability}
\begin{assignmentinstruction}
Determine the conditions for a stable equilibrium of your unforced ($\mu=0$) and damped autonomous subsystem. Formulate the range of parameters that ensure their stability [$\operatorname{Re}(\lambda_n)<0$, $n=1$--$4$].
\end{assignmentinstruction}

\subsubsection*{Equilibrium conditions}
For an autonomous equilibrium, set the magnetic forcing to zero, so that $\mu=0$ and $f_1=0$, and require all velocities and accelerations to vanish. Writing $(y_1,y_2)=(y_{1e},y_{2e})$, the equilibrium must lie in the physical domain
\[
1-(y_{1e}\pm y_{2e})>0,
\qquad\text{equivalently}\qquad y_{1e}+|y_{2e}|<1.
\]
The resulting coupled algebraic equations are
\begin{align}
\omega_1^2 y_{1e}+\alpha_{11}y_{1e}^3+\alpha_{12}y_{1e}y_{2e}^2
&=\beta\left[\frac{1}{(1-y_{1e}-y_{2e})^2}+\frac{1}{(1-y_{1e}+y_{2e})^2}\right],\label{eq:eq-y1}\\
\omega_2^2 y_{2e}+\alpha_{21}y_{1e}^2y_{2e}+\alpha_{22}y_{2e}^3
&=\beta\left[\frac{1}{(1-y_{1e}-y_{2e})^2}-\frac{1}{(1-y_{1e}+y_{2e})^2}\right].\label{eq:eq-y2}
\end{align}
Because the second coordinate represents the antisymmetric mode, setting $y_{2e}=0$ makes both sides of \eqref{eq:eq-y2} vanish. This gives the fundamental symmetric branch $(y_{1e},0)$, with $0<y_{1e}<1$ determined by
\begin{equation}\label{eq:symmetric-equilibrium}
\omega_1^2 y_{1e}+\alpha_{11}y_{1e}^3=\frac{2\beta}{(1-y_{1e})^2}.
\end{equation}
Possible asymmetric branches with $y_{2e}\ne0$ must additionally satisfy \eqref{eq:eq-y1} and, after division of \eqref{eq:eq-y2} by $y_{2e}$, the corresponding nonzero-mode balance. Such branches can arise if electrostatic softening overcomes the antisymmetric mechanical restoring stiffness.

\subsubsection*{Linear stability conditions}
Introduce perturbations $x_1,x_2$ about an equilibrium by $y_i=y_{ie}+x_i$. To first order, the perturbations satisfy
\begin{equation}\label{eq:linearized-system}
\mathbf{x}_{\tau\tau}+\mathbf{C}\mathbf{x}_{\tau}+\mathbf{K}\mathbf{x}=\mathbf{0},
\qquad \mathbf{x}=(x_1,x_2)^T,
\end{equation}
where the tangent stiffness entries, evaluated at a general equilibrium, are
\begin{align*}
K_{11}&=\omega_1^2+3\alpha_{11}y_{1e}^2+\alpha_{12}y_{2e}^2
-2\beta\left[\frac{1}{(1-y_{1e}-y_{2e})^3}+\frac{1}{(1-y_{1e}+y_{2e})^3}\right],\\
K_{12}&=2\alpha_{12}y_{1e}y_{2e}
-2\beta\left[\frac{1}{(1-y_{1e}-y_{2e})^3}-\frac{1}{(1-y_{1e}+y_{2e})^3}\right],\\
K_{21}&=2\alpha_{21}y_{1e}y_{2e}
-2\beta\left[\frac{1}{(1-y_{1e}-y_{2e})^3}-\frac{1}{(1-y_{1e}+y_{2e})^3}\right],\\
K_{22}&=\omega_2^2+\alpha_{21}y_{1e}^2+3\alpha_{22}y_{2e}^2
-2\beta\left[\frac{1}{(1-y_{1e}-y_{2e})^3}+\frac{1}{(1-y_{1e}+y_{2e})^3}\right].
\end{align*}
The linearized damping matrix follows from retaining the terms linear in perturbation velocities:
\[
\mathbf{C}=\begin{pmatrix}
\zeta+\eta_{11}y_{1e}^2 & \eta_{12}y_{1e}y_{2e}\\
\eta_{11}y_{1e}^2 & \zeta+\eta_{12}y_{1e}y_{2e}
\end{pmatrix}.
\]
In particular, at the symmetric equilibrium, define $\hat{\zeta}_1=\zeta+\eta_{11}y_{1e}^2$. The $\eta_{11}$ term contributes to the damping of $x_1$ and also gives a one-way velocity coupling from $x_{1,\tau}$ into the $x_2$ equation; the $\eta_{12}$ contribution vanishes there. Reflection symmetry makes the tangent stiffness diagonal:
\begin{equation}\label{eq:symmetric-stiffness}
\mathbf{K}=\begin{pmatrix}K_{11}&0\\0&K_{22}\end{pmatrix},\qquad
K_{11}=\omega_1^2+3\alpha_{11}y_{1e}^2-\frac{4\beta}{(1-y_{1e})^3},\quad
K_{22}=\omega_2^2+\alpha_{21}y_{1e}^2-\frac{4\beta}{(1-y_{1e})^3}.
\end{equation}
For exponential perturbations $\mathbf{x}=\mathbf{v}e^{\lambda\tau}$, the characteristic equation is $P(\lambda)=\det(\lambda^2\mathbf{I}+\lambda\mathbf{C}+\mathbf{K})=0$. Although $\mathbf{C}$ is lower triangular rather than diagonal, the stiffness matrix is diagonal on this branch, and hence
\begin{equation}\label{eq:symmetric-characteristic}
P(\lambda)=(\lambda^2+\hat{\zeta}_1\lambda+K_{11})(\lambda^2+\zeta\lambda+K_{22})=0.
\end{equation}
Each quadratic has both roots in the open left half-plane precisely when its damping and stiffness coefficients are positive. Thus, for the symmetric branch, asymptotic stability requires
\begin{equation}\label{eq:stability-conditions}
\zeta>0,\qquad \hat{\zeta}_1=\zeta+\eta_{11}y_{1e}^2>0,\qquad K_{11}>0,\qquad K_{22}>0.
\end{equation}
With the stipulated positive damping parameters, the damping inequalities hold. The two stiffness inequalities are equivalently the bound
\begin{equation}\label{eq:beta-stability-bound}
\beta<\frac{(1-y_{1e})^3}{4}\min\!\left(\omega_1^2+3\alpha_{11}y_{1e}^2,\;\omega_2^2+\alpha_{21}y_{1e}^2\right).
\end{equation}
This is an implicit parameter restriction because $y_{1e}$ itself depends on $\beta$ through \eqref{eq:symmetric-equilibrium}. Loss of $K_{11}>0$ is the symmetric-mode pull-in instability and, generically, marks a fold of the equilibrium branch; loss of $K_{22}>0$ destabilizes the antisymmetric mode and can produce a symmetry-breaking (pitchfork) bifurcation. If nonlinear damping is instead assigned to higher perturbation order in the subsequent asymptotic analysis, the leading linearization has $\mathbf{C}=\zeta\mathbf{I}$; for $\zeta>0$ the same stiffness conditions determine stability, while retaining the damping terms gives the more precise eigenvalues in \eqref{eq:symmetric-characteristic}.

\subsection{Natural Frequencies and Internal Resonance}
\begin{assignmentinstruction}
Derive the undamped! natural frequencies of the linearized dynamical system about the stable equilibrium [e.g. $\omega_{1,2}$ from $\mathbf{A}\mathbf{y}_{\tau\tau}+\mathbf{B}\mathbf{y}=0$ where $\mathbf{y}=(y_1,y_2)^T$ and $\mathbf{A}$, $\mathbf{B}$ are constants].

Note possible conditions for an internal resonance [$\omega_2\sim m\omega_1$ where $m=1,2,3$].
\end{assignmentinstruction}

Set $\zeta=\eta_{ij}=0$ and $\mu=0$, and consider small displacements about the stable symmetric equilibrium. The undamped linearized equations take the requested form
\begin{equation}\label{eq:undamped-eigenproblem}
\mathbf{A}\mathbf{x}_{\tau\tau}+\mathbf{B}\mathbf{x}=\mathbf{0},\qquad
\mathbf{x}=(x_1,x_2)^T,\quad \mathbf{A}=\mathbf{I}_2,\quad
\mathbf{B}=\mathbf{K}=\operatorname{diag}(\hat{\omega}_1^2,\hat{\omega}_2^2).
\end{equation}
Consequently, the equilibrium-shifted natural frequencies are
\begin{align}
\hat{\omega}_1&=\sqrt{\omega_1^2+3\alpha_{11}y_{1e}^2-\frac{4\beta}{(1-y_{1e})^3}},\label{eq:frequency-one}\\
\hat{\omega}_2&=\sqrt{\omega_2^2+\alpha_{21}y_{1e}^2-\frac{4\beta}{(1-y_{1e})^3}}.\label{eq:frequency-two}
\end{align}
The square roots are real on the linearly stable branch. For a nonsymmetric equilibrium, $\mathbf{B}=\mathbf{K}$ is generally not diagonal; its squared natural frequencies are the eigenvalues obtained from $\det(\mathbf{B}-\hat{\omega}^{\,2}\mathbf{A})=0$. In later perturbation calculations, the shifted frequencies may be denoted simply by $\omega_1,\omega_2$ when no confusion with the bare parameters arises.

The electrostatic load tunes the frequencies through both the equilibrium shift $y_{1e}(\beta)$ and the negative tangent-stiffness contribution $-4\beta/(1-y_{1e})^3$. This term is the electrostatic softening contribution; the total frequency variation with $\beta$ also reflects the change in $y_{1e}$ and need not be identical for the two modes, since the mechanical stiffness contributions $3\alpha_{11}y_{1e}^2$ and $\alpha_{21}y_{1e}^2$ generally differ and the bare frequencies may be unequal.

An $m$-to-one internal resonance, with $m\in\{1,2,3\}$, requires the frequency mismatch to be small on the perturbation scale, equivalently
\begin{equation}\label{eq:internal-resonance-general}
\hat{\omega}_2\approx m\hat{\omega}_1
\quad\Longleftrightarrow\quad
\omega_2^2+\alpha_{21}y_{1e}^2-\frac{4\beta}{(1-y_{1e})^3}
\approx m^2\left(\omega_1^2+3\alpha_{11}y_{1e}^2-\frac{4\beta}{(1-y_{1e})^3}\right),
\end{equation}
where the squared-frequency mismatch is $\mathcal{O}(\varepsilon)$ in a multiple-scales formulation. The three cases reduce to
\begin{align*}
m=1:\quad &\omega_2^2-\omega_1^2\approx(3\alpha_{11}-\alpha_{21})y_{1e}^2,\\
m=2:\quad &\omega_2^2-4\omega_1^2\approx(12\alpha_{11}-\alpha_{21})y_{1e}^2-\frac{12\beta}{(1-y_{1e})^3},\\
m=3:\quad &\omega_2^2-9\omega_1^2\approx(27\alpha_{11}-\alpha_{21})y_{1e}^2-\frac{32\beta}{(1-y_{1e})^3}.
\end{align*}
For $1{:}1$ resonance the common electrostatic tangent-stiffness term cancels exactly. The $1{:}2$ condition is associated with quadratic interaction terms, which arise upon expanding the electrostatic forces about the generally nonzero equilibrium; the $1{:}3$ condition is associated with cubic interactions from the mechanical and electrostatic nonlinearities. These are frequency-matching conditions; whether a given resonance generates appreciable modal energy transfer also depends on the corresponding nonlinear coupling coefficients and symmetry selection rules.

\section{Asymptotic Analysis of the Modulated System}
\subsection{Expansion and Perturbation Ordering}
\begin{assignmentinstruction}
Consider small displacements about the stable equilibrium position and expand the equations of motion above to cubic order where the damping and forcing parameters are sufficiently small. Rewrite your equations of motion in the form ready singular perturbation.
\[
\begin{aligned}
x_{1,\tau\tau}+\omega_1^2x_1 &= \varepsilon F_1(x_1,x_{1,\tau},x_2,x_{2,\tau},\tau),\\
x_{2,\tau\tau}+\omega_2^2x_2 &= \varepsilon F_2(x_1,x_{1,\tau},x_2,x_{2,\tau},\tau).
\end{aligned}
\]
Where the right-hand-sides include: (i) linear damping terms, (ii) forcing term, (iii) quadratic and cubic contributions of the interaction terms.
\end{assignmentinstruction}

Set $d=1-y_{1e}>0$ and write $y_1=y_{1e}+x_1$, $y_2=x_2$. For $|x_1\pm x_2|<d$, the electrostatic expansion is
\begin{equation}
\frac{1}{[1-(y_{1e}+x_1\pm x_2)]^2}
=\frac{1}{d^2}\left[1+\frac{2(x_1\pm x_2)}{d}+\frac{3(x_1\pm x_2)^2}{d^2}+\frac{4(x_1\pm x_2)^3}{d^3}\right]+\mathcal{O}(\|\mathbf{x}\|^4),
\end{equation}
and hence
\begin{align*}
g_1&=\frac{2\beta}{d^2}+\frac{4\beta}{d^3}x_1+\frac{6\beta}{d^4}(x_1^2+x_2^2)+\frac{8\beta}{d^5}(x_1^3+3x_1x_2^2)+\mathcal{O}(\|\mathbf{x}\|^4),\\
g_2&=\frac{4\beta}{d^3}x_2+\frac{12\beta}{d^4}x_1x_2+\frac{8\beta}{d^5}(3x_1^2x_2+x_2^3)+\mathcal{O}(\|\mathbf{x}\|^4).
\end{align*}
The mechanical terms expand as
\begin{align*}
\alpha_{11}y_1^3&=\alpha_{11}y_{1e}^3+3\alpha_{11}y_{1e}^2x_1+3\alpha_{11}y_{1e}x_1^2+\alpha_{11}x_1^3,\\
\alpha_{12}y_1y_2^2&=\alpha_{12}y_{1e}x_2^2+\alpha_{12}x_1x_2^2,\\
\alpha_{21}y_1^2y_2&=\alpha_{21}y_{1e}^2x_2+2\alpha_{21}y_{1e}x_1x_2+\alpha_{21}x_1^2x_2,\\
\alpha_{22}y_2^3&=\alpha_{22}x_2^3,
\end{align*}
while the nonlinear damping terms have the expansions
\begin{align*}
\eta_{11}y_1^2x_{1,\tau}&=\eta_{11}y_{1e}^2x_{1,\tau}+2\eta_{11}y_{1e}x_1x_{1,\tau}+\eta_{11}x_1^2x_{1,\tau},\\
\eta_{12}y_1y_2x_{2,\tau}&=\eta_{12}y_{1e}x_2x_{2,\tau}+\eta_{12}x_1x_2x_{2,\tau}.
\end{align*}
After subtracting the equilibrium balance and moving the tangent-linear terms to the left, use the coefficients $\hat{\omega}_1,\hat{\omega}_2$ from Section 2 and define
\begin{gather*}
\hat{\zeta}_1=\zeta+\eta_{11}y_{1e}^2,\quad \hat{\zeta}_2=\zeta,\quad c_{21}=\eta_{11}y_{1e}^2,\\
q_{11}=\frac{6\beta}{d^4}-3\alpha_{11}y_{1e},\quad q_{12}=\frac{6\beta}{d^4}-\alpha_{12}y_{1e},\quad q_{21}=\frac{12\beta}{d^4}-2\alpha_{21}y_{1e},\\
\kappa_{11}=\alpha_{11}-\frac{8\beta}{d^5},\quad \kappa_{12}=\alpha_{12}-\frac{24\beta}{d^5},\quad
\kappa_{21}=\alpha_{21}-\frac{24\beta}{d^5},\quad \kappa_{22}=\alpha_{22}-\frac{8\beta}{d^5},\\
\tilde{\eta}_{11}=2\eta_{11}y_{1e},\qquad \tilde{\eta}_{12}=\eta_{12}y_{1e}.
\end{gather*}
The resulting truncated equations, with a formal factor $\varepsilon$ marking weak damping, forcing, and nonlinear corrections, are
\begin{align}
x_{1,\tau\tau}+\hat{\omega}_1^2x_1&=\varepsilon F_1,\nonumber\\
\varepsilon F_1&=\mu\sin(\Omega\tau)-\hat{\zeta}_1x_{1,\tau}-\tilde{\eta}_{11}x_1x_{1,\tau}-\eta_{11}x_1^2x_{1,\tau}-\tilde{\eta}_{12}x_2x_{2,\tau}-\eta_{12}x_1x_2x_{2,\tau}\nonumber\\
&\hspace{1.4em}+q_{11}x_1^2+q_{12}x_2^2-\kappa_{11}x_1^3-\kappa_{12}x_1x_2^2+\mathcal{O}(\|\mathbf{x}\|^4),\label{eq:F1-expanded}\\
x_{2,\tau\tau}+\hat{\omega}_2^2x_2&=\varepsilon F_2,\nonumber\\
\varepsilon F_2&=-\hat{\zeta}_2x_{2,\tau}-c_{21}x_{1,\tau}-\tilde{\eta}_{11}x_1x_{1,\tau}-\eta_{11}x_1^2x_{1,\tau}-\tilde{\eta}_{12}x_2x_{2,\tau}-\eta_{12}x_1x_2x_{2,\tau}\nonumber\\
&\hspace{1.4em}+q_{21}x_1x_2-\kappa_{21}x_1^2x_2-\kappa_{22}x_2^3+\mathcal{O}(\|\mathbf{x}\|^4).\label{eq:F2-expanded}
\end{align}
Here the displayed left-hand frequencies are the equilibrium-shifted frequencies $\hat{\omega}_{1,2}$ (the $\omega_{1,2}$ in the subsequent assignment template). The perturbation ordering means the damping coefficients, $\mu$, and the retained nonlinear corrections are formally $\mathcal{O}(\varepsilon)$; equivalently, one may regard their coefficients in $F_k$ as the rescaled order-one quantities after extracting $\varepsilon$. The model as written contains the cross-velocity term $c_{21}x_{1,\tau}$ and the $\eta_{11}$ damping terms in the second equation; these are retained rather than imposing an additional modal-parity assumption.

\subsection{Order Equations}
\begin{assignmentinstruction}
Derive the order equations where the first order includes only the linear conservative components. Note your small parameters clearly.
\end{assignmentinstruction}
Introduce $T_n=\varepsilon^n\tau$, with $T_0=\tau$, $T_1=\varepsilon\tau$, and $D_n=\partial/\partial T_n$. Then
\[
\frac{d}{d\tau}=D_0+\varepsilon D_1+\mathcal{O}(\varepsilon^2),\qquad
\frac{d^2}{d\tau^2}=D_0^2+2\varepsilon D_0D_1+\mathcal{O}(\varepsilon^2),
\]
and expand $x_k=x_{k0}(T_0,T_1)+\varepsilon x_{k1}(T_0,T_1)+\mathcal{O}(\varepsilon^2)$ for $k=1,2$. The weak-damping/forcing ordering is $\zeta,\eta_{ij},\mu=\mathcal{O}(\varepsilon)$ in the unscaled equations; the nonlinear terms retained in \eqref{eq:F1-expanded}--\eqref{eq:F2-expanded} are counted at the same perturbation order (for example, by a weak-nonlinearity or small-amplitude ordering).
\begin{align}
\mathcal{O}(1):\qquad&D_0^2x_{10}+\hat{\omega}_1^2x_{10}=0,&
D_0^2x_{20}+\hat{\omega}_2^2x_{20}=0,\label{eq:order-zero}\\[2pt]
\mathcal{O}(\varepsilon):\qquad&D_0^2x_{11}+\hat{\omega}_1^2x_{11}=-2D_0D_1x_{10}+F_1^{(0)},&
D_0^2x_{21}+\hat{\omega}_2^2x_{21}=-2D_0D_1x_{20}+F_2^{(0)},\label{eq:order-one}
\end{align}
where, with all variables and velocities evaluated at leading order,
\begin{align*}
F_1^{(0)}={}&\mu\sin(\Omega T_0)-\hat{\zeta}_1D_0x_{10}-\tilde{\eta}_{11}x_{10}D_0x_{10}-\eta_{11}x_{10}^2D_0x_{10}-\tilde{\eta}_{12}x_{20}D_0x_{20}-\eta_{12}x_{10}x_{20}D_0x_{20}\\
&+q_{11}x_{10}^2+q_{12}x_{20}^2-\kappa_{11}x_{10}^3-\kappa_{12}x_{10}x_{20}^2,\\
F_2^{(0)}={}&-\hat{\zeta}_2D_0x_{20}-c_{21}D_0x_{10}-\tilde{\eta}_{11}x_{10}D_0x_{10}-\eta_{11}x_{10}^2D_0x_{10}-\tilde{\eta}_{12}x_{20}D_0x_{20}-\eta_{12}x_{10}x_{20}D_0x_{20}\\
&+q_{21}x_{10}x_{20}-\kappa_{21}x_{10}^2x_{20}-\kappa_{22}x_{20}^3.
\end{align*}

\subsection{Solvability Conditions}
\begin{assignmentinstruction}
Assume a periodic multiple scale solution form for the homogeneous first order equations and derive the solvability conditions required to ensure that the asymptotic solution does NOT grow without bound.
\end{assignmentinstruction}
The leading-order periodic solutions are
\[
x_{10}=A_1(T_1)e^{i\hat{\omega}_1T_0}+\bar A_1(T_1)e^{-i\hat{\omega}_1T_0},\qquad
x_{20}=A_2(T_1)e^{i\hat{\omega}_2T_0}+\bar A_2(T_1)e^{-i\hat{\omega}_2T_0}.
\]
For primary external resonance and internal resonance, introduce
\[
\Omega=\hat{\omega}_1+\varepsilon\sigma_1,\qquad
\hat{\omega}_2=m\hat{\omega}_1+\varepsilon\sigma_2,\quad m\in\{1,2,3\},
\]
so that $\mu\sin(\Omega T_0)=-\frac{i\mu}{2}e^{i\sigma_1T_1}e^{i\hat{\omega}_1T_0}+\frac{i\mu}{2}e^{-i\sigma_1T_1}e^{-i\hat{\omega}_1T_0}$. Secular terms are the coefficients of $e^{i\hat{\omega}_kT_0}$ in the corresponding right-hand sides of \eqref{eq:order-one}. The universal terms are
\begin{align*}
S_1^{(0)}={}&-2i\hat{\omega}_1D_1A_1-i\hat{\omega}_1\hat{\zeta}_1A_1
-(3\kappa_{11}+i\hat{\omega}_1\eta_{11})|A_1|^2A_1-2\kappa_{12}|A_2|^2A_1
-\frac{i\mu}{2}e^{i\sigma_1T_1},\\
S_2^{(0)}={}&-2i\hat{\omega}_2D_1A_2-i\hat{\omega}_2\hat{\zeta}_2A_2
-3\kappa_{22}|A_2|^2A_2-2\kappa_{21}|A_1|^2A_2.
\end{align*}
The first line includes the resonant component of $-\eta_{11}x_{10}^2D_0x_{10}$. In the second line the last damping term sometimes associated with $|A_1|^2A_2$ is absent: in the equations stated in Section 1, the $\eta_{11}$ term in the second equation is $-\eta_{11}x_{10}^2D_0x_{10}$, not $-\eta_{11}x_{10}^2D_0x_{20}$. Its resonant contributions instead occur at $\hat{\omega}_2=\hat{\omega}_1$, $2\hat{\omega}_1$, or $3\hat{\omega}_1$ as described below.

For $1{:}1$ resonance ($m=1$), the additional resonant terms are
\begin{align*}
S_1^{(1)}={}&-(\kappa_{12}+i\hat{\omega}_2\eta_{12})A_2^2\bar A_1e^{2i\sigma_2T_1},\\
S_2^{(1)}={}&-i\hat{\omega}_1c_{21}A_1e^{-i\sigma_2T_1}
-i\hat{\omega}_1\eta_{11}|A_1|^2A_1e^{-i\sigma_2T_1}
-\kappa_{21}A_1^2\bar A_2e^{-2i\sigma_2T_1}
-i\hat{\omega}_2\eta_{12}\bar A_1A_2^2e^{i\sigma_2T_1}.
\end{align*}
The first and third terms are the phase-sensitive cubic stiffness couplings; the $\eta_{12}$ terms are their velocity-dependent counterparts. The $c_{21}$ term is a direct linear cross-velocity excitation of mode 2. The $\eta_{11}$ term shown in $S_2^{(1)}$ follows from the damping term as it is actually written in the model.

For $1{:}2$ resonance ($\hat{\omega}_2=2\hat{\omega}_1+\varepsilon\sigma_2$), the conservative quadratic coupling permitted on a general asymmetric or imperfect branch has the form
\[
S_1^{(2)}=\Gamma_{12}\bar A_1A_2e^{i\sigma_2T_1},\qquad
S_2^{(2)}=\left(\Gamma_{21}-i\hat{\omega}_1\tilde{\eta}_{11}\right)A_1^2e^{-i\sigma_2T_1}.
\]
Here $\Gamma_{12}$ is the coefficient of $x_1x_2$ in $F_1$, projected at $\hat{\omega}_2-\hat{\omega}_1$, and $\Gamma_{21}$ is the coefficient of $x_1^2$ in $F_2$, projected at $2\hat{\omega}_1$ (including any velocity-dependent quadratic contribution). On the exact symmetric branch, the conservative $x_1x_2$ and $x_1^2$ stiffness terms vanish by reflection parity: the conservative $F_1$ is even in $x_2$, whereas conservative $F_2$ is odd in $x_2$. For the equations given here, however, the second-equation damping term breaks that parity and contributes
\[
S_{2,\mathrm{quad}}^{(\eta)}=-i\hat{\omega}_1\tilde{\eta}_{11}A_1^2e^{-i\sigma_2T_1}.
\]
Thus $S_1=S_1^{(0)}$ and $S_2=S_2^{(0)}+S_2^{(2)}$ on the perfectly symmetric branch when $\Gamma_{12}=\Gamma_{21}=0$.

For the reverse, symmetric-branch $2{:}1$ resonance, let $\hat{\omega}_1=2\hat{\omega}_2+\varepsilon\tilde{\sigma}_2$ and $\tilde{\gamma}_2=\theta_1-2\theta_2+\tilde{\sigma}_2T_1$. The retained quadratic terms in \eqref{eq:F1-expanded}--\eqref{eq:F2-expanded} give
\begin{equation}\label{eq:solvability-21}
S_1^{(2:1)}=(q_{12}-i\hat{\omega}_2\tilde{\eta}_{12})A_2^2e^{-i\tilde{\sigma}_2T_1},\qquad
S_2^{(2:1)}=q_{21}A_1\bar A_2e^{i\tilde{\sigma}_2T_1}.
\end{equation}
Unlike the $1{:}2$ coupling in the conservative equations, this reverse $2{:}1$ interaction is allowed on the symmetric branch: both $q_{12}$ and $q_{21}$ are generally nonzero, and the first-mode equation also contains the velocity-dependent coefficient $\tilde{\eta}_{12}$. It therefore provides bidirectional quadratic modal coupling without symmetry breaking.

For $1{:}3$ resonance ($\hat{\omega}_2=3\hat{\omega}_1+\varepsilon\sigma_2$), generic symmetry-breaking cubic stiffness terms give
\[
S_1^{(3)}=\Lambda_{12}\bar A_1^{\,2}A_2e^{i\sigma_2T_1},\qquad
S_2^{(3)}=(\Lambda_{21}-i\hat{\omega}_1\eta_{11})A_1^3e^{-i\sigma_2T_1}.
\]
These conservative coefficients vanish for the relevant $x_1^2x_2$ term in $F_1$ and $x_1^3$ term in $F_2$ on the symmetric branch. The specified second-equation nonlinear damping nevertheless supplies the resonant term
\[
S_{2,\mathrm{cub}}^{(\eta)}=-i\hat{\omega}_1\eta_{11}A_1^3e^{-i\sigma_2T_1}.
\]
On the symmetric branch, the remaining cubic cross-modal effects are the cross-frequency (cross-Kerr) terms in $S_1^{(0)}$ and $S_2^{(0)}$, the $1{:}1$ phase-coupling terms above, and the damping terms retained explicitly. Finally, the complex solvability conditions are
\begin{equation}\label{eq:solvability}
S_1^{(0)}+S_1^{(m)}=0,\qquad S_2^{(0)}+S_2^{(m)}=0,
\end{equation}
where for each resonance $m$, the additional terms are those just stated; these conditions remove all resonant forcing and prevent secular growth.

\subsection{Slow-Flow Equations}
\begin{assignmentinstruction}
Derive the coupled slowly varying evolution equations from the (complex) solvability condition.
\end{assignmentinstruction}
Write $A_k=\tfrac12a_ke^{i\theta_k}$, $a_k\ge0$, and define $\gamma_1=\sigma_1T_1-\theta_1$. With the detuning convention $\hat{\omega}_2=m\hat{\omega}_1+\varepsilon\sigma_2$ used above, the autonomous internal phase is
\begin{equation}\label{eq:relative-phase}
\gamma_2=m\theta_1-\theta_2-\sigma_2T_1.
\end{equation}
The minus sign is required by this detuning convention: a resonant product $A_1^m e^{-i\sigma_2T_1}$ has phase $m\theta_1-\theta_2-\sigma_2T_1$. If the alternative convention $\hat{\omega}_2=m\hat{\omega}_1-\varepsilon\sigma_2$ is adopted, the plus sign in the relative-phase definition is recovered.

Below, all primes denote $d/dT_1$. Separating real and imaginary parts of the solvability equations gives explicit autonomous systems. For the $1{:}2$ case, $\hat{\omega}_2=2\hat{\omega}_1+\varepsilon\sigma_2$, with $\gamma_2=2\theta_1-\theta_2-\sigma_2T_1$, the general quadratic-coupling flow is
\begin{align}
a_1'={}&-\frac{\hat{\zeta}_1}{2}a_1-\frac{\eta_{11}}{8}a_1^3-\frac{\Gamma_{12}}{4\hat{\omega}_1}a_1a_2\sin\gamma_2-\frac{\mu}{2\hat{\omega}_1}\cos\gamma_1,\nonumber\\
a_1\gamma_1'={}&\sigma_1a_1-\frac{3\kappa_{11}}{8\hat{\omega}_1}a_1^3-\frac{\kappa_{12}}{4\hat{\omega}_1}a_1a_2^2+\frac{\Gamma_{12}}{4\hat{\omega}_1}a_1a_2\cos\gamma_2+\frac{\mu}{2\hat{\omega}_1}\sin\gamma_1,\nonumber\\
a_2'={}&-\frac{\hat{\zeta}_2}{2}a_2+\frac{a_1^2}{4\hat{\omega}_2}(\Gamma_{21}\sin\gamma_2-\hat{\omega}_1\tilde{\eta}_{11}\cos\gamma_2),\nonumber\\
a_2\gamma_2'={}&-\sigma_2a_2+2a_2(\sigma_1-\gamma_1')-\frac{3\kappa_{22}}{8\hat{\omega}_2}a_2^3-\frac{\kappa_{21}}{4\hat{\omega}_2}a_1^2a_2
\frac{a_1^2}{4\hat{\omega}_2}(\Gamma_{21}\cos\gamma_2+\hat{\omega}_1\tilde{\eta}_{11}\sin\gamma_2).\label{eq:slowflow-12}
\end{align}
On the exact symmetric branch the conservative $\Gamma_{12},\Gamma_{21}$ vanish; the displayed $\tilde{\eta}_{11}$ terms remain if the nonlinear damping in the second equation is retained.

For the exact symmetric-branch reverse $2{:}1$ resonance, $\hat{\omega}_1=2\hat{\omega}_2+\varepsilon\tilde{\sigma}_2$ and $\tilde{\gamma}_2=\theta_1-2\theta_2+\tilde{\sigma}_2T_1$. The nonzero quadratic coefficients $q_{12},q_{21}$ give
\begin{align}
a_1'={}&-\frac{\hat{\zeta}_1}{2}a_1-\frac{\eta_{11}}{8}a_1^3-\frac{a_2^2}{4\hat{\omega}_1}(q_{12}\sin\tilde{\gamma}_2+\hat{\omega}_2\tilde{\eta}_{12}\cos\tilde{\gamma}_2)-\frac{\mu}{2\hat{\omega}_1}\cos\gamma_1,\nonumber\\
a_1\gamma_1'={}&\sigma_1a_1-\frac{3\kappa_{11}}{8\hat{\omega}_1}a_1^3-\frac{\kappa_{12}}{4\hat{\omega}_1}a_1a_2^2
+\frac{a_2^2}{4\hat{\omega}_1}(q_{12}\cos\tilde{\gamma}_2-\hat{\omega}_2\tilde{\eta}_{12}\sin\tilde{\gamma}_2)+\frac{\mu}{2\hat{\omega}_1}\sin\gamma_1,\nonumber\\
a_2'={}&-\frac{\hat{\zeta}_2}{2}a_2+\frac{q_{21}}{4\hat{\omega}_2}a_1a_2\sin\tilde{\gamma}_2,\nonumber\\
a_2\tilde{\gamma}_2'={}&(\tilde{\sigma}_2+\sigma_1-\gamma_1')a_2-\frac{3\kappa_{22}}{4\hat{\omega}_2}a_2^3-\frac{\kappa_{21}}{2\hat{\omega}_2}a_1^2a_2
+\frac{q_{21}}{2\hat{\omega}_2}a_1a_2\cos\tilde{\gamma}_2.\label{eq:slowflow-21}
\end{align}
The $q_{12}$ term in mode 1 and $q_{21}$ term in mode 2 provide the bidirectional conservative interaction; the $\tilde{\eta}_{12}$ term adds velocity-dependent coupling. This quadratic interaction is permitted by the symmetric-branch mode parity.

For $1{:}3$ resonance, $\hat{\omega}_2=3\hat{\omega}_1+\varepsilon\sigma_2$ and $\gamma_2=3\theta_1-\theta_2-\sigma_2T_1$, the explicit cubic flow is
\begin{align}
a_1'={}&-\frac{\hat{\zeta}_1}{2}a_1-\frac{\eta_{11}}{8}a_1^3-\frac{\Lambda_{12}}{8\hat{\omega}_1}a_1^2a_2\sin\gamma_2-\frac{\mu}{2\hat{\omega}_1}\cos\gamma_1,\nonumber\\
a_1\gamma_1'={}&\sigma_1a_1-\frac{3\kappa_{11}}{8\hat{\omega}_1}a_1^3-\frac{\kappa_{12}}{4\hat{\omega}_1}a_1a_2^2+\frac{\Lambda_{12}}{8\hat{\omega}_1}a_1^2a_2\cos\gamma_2+\frac{\mu}{2\hat{\omega}_1}\sin\gamma_1,\nonumber\\
a_2'={}&-\frac{\hat{\zeta}_2}{2}a_2+\frac{a_1^3}{8\hat{\omega}_2}(\Lambda_{21}\sin\gamma_2-\hat{\omega}_1\eta_{11}\cos\gamma_2),\nonumber\\
a_2\gamma_2'={}&-\sigma_2a_2+3a_2(\sigma_1-\gamma_1')-\frac{3\kappa_{22}}{8\hat{\omega}_2}a_2^3-\frac{\kappa_{21}}{4\hat{\omega}_2}a_1^2a_2
+\frac{a_1^3}{8\hat{\omega}_2}(\Lambda_{21}\cos\gamma_2+\hat{\omega}_1\eta_{11}\sin\gamma_2).\label{eq:slowflow-13}
\end{align}
For $1{:}1$ resonance, set $\gamma_2=\theta_1-\theta_2-\sigma_2T_1$. The terms added to the universal flow are explicitly
\begin{align*}
\delta a_1'={}&\frac{a_1a_2^2}{8\hat{\omega}_1}(\kappa_{12}\sin2\gamma_2-\hat{\omega}_2\eta_{12}\cos2\gamma_2),\\
\delta(a_1\gamma_1')={}&-\frac{a_2^2}{8\hat{\omega}_1}(\kappa_{12}\cos2\gamma_2+\hat{\omega}_2\eta_{12}\sin2\gamma_2),\\
\delta a_2'={}&-\frac{\hat{\omega}_1a_1}{2\hat{\omega}_2}c_{21}\cos\gamma_2-\frac{\hat{\omega}_1\eta_{11}}{8\hat{\omega}_2}a_1^3\cos\gamma_2-\frac{\kappa_{21}}{8\hat{\omega}_2}a_1^2a_2\sin2\gamma_2-\frac{\eta_{12}}{8}a_1a_2^2\cos\gamma_2,\\
\delta(a_2\gamma_2')={}&-a_2\,\delta\gamma_1'+\frac{\hat{\omega}_1c_{21}}{2\hat{\omega}_2}a_1\sin\gamma_2+\frac{\hat{\omega}_1\eta_{11}}{8\hat{\omega}_2}a_1^3\sin\gamma_2-\frac{\kappa_{21}}{8\hat{\omega}_2}a_1^2a_2\cos2\gamma_2-\frac{\eta_{12}}{8}a_1a_2^2\sin\gamma_2,
\end{align*}
where $\delta\gamma_1'=\delta(a_1\gamma_1')/a_1$. In all cases, the nonresonant single-mode limit is $a_2=0$, with the primary-resonance equations obtained by dropping the coupling terms.

\subsection{Amplitude Response and Stability}
\begin{assignmentinstruction}
Determine the non-trivial equilibria of the amplitude evolution equations (which correspond to the amplitudes of the gyro dependent variables) and formulate the conditions of stability of these equilibria.
\end{assignmentinstruction}
\subsubsection{Steady Amplitudes}
Steady periodic responses correspond to fixed points of the slow flow, $a_1'=\gamma_1'=a_2'=\gamma_2'=0$, or equivalently to the complex algebraic equations $G_1=G_2=0$. For an uncoupled or nonresonant single-mode response ($a_2=0$), the first-mode equations give
\begin{equation}\label{eq:single-mode-response}
\left[\frac{\hat{\zeta}_1}{2}a_{10}+\frac{\eta_{11}}{8}a_{10}^3\right]^2
+\left[\sigma_1a_{10}-\frac{3\kappa_{11}}{8\hat{\omega}_1}a_{10}^3\right]^2
=\left(\frac{\mu}{2\hat{\omega}_1}\right)^2.
\end{equation}
This is the nonlinearly damped Duffing response curve. The assignment's $a_{20}=0$ branch is an exact solution only when the resonant forcing of mode 2 vanishes. For the model as printed, $c_{21}x_{1,\tau}$ forces mode 2 at $1{:}1$, $\tilde{\eta}_{11}x_1x_{1,\tau}$ supplies a $1{:}2$ resonant term, and $\eta_{11}x_1^2x_{1,\tau}$ supplies a $1{:}3$ resonant term. Thus, at those internal resonances, the single-mode curve is an approximation unless the corresponding damping coupling is negligible or suppressed.

For each coupled steady response, set the four right-hand sides of the appropriate real slow-flow system to zero. The $1{:}2$ branch ($\hat{\omega}_2=2\hat{\omega}_1+\varepsilon\sigma_2$) can be written without an unevaluated complex modulus as follows. Define
\begin{align*}
U_2&=\frac{2\hat{\omega}_2\hat{\zeta}_2a_{20}}{a_{10}^2},\\
V_2&=\frac{4\hat{\omega}_2a_{20}}{a_{10}^2}\left(\sigma_2-2\sigma_1+\frac{3\kappa_{22}}{8\hat{\omega}_2}a_{20}^2+\frac{\kappa_{21}}{4\hat{\omega}_2}a_{10}^2\right),\qquad d_2=\hat{\omega}_1\tilde{\eta}_{11},\\
\sin\gamma_{20}&=\frac{\Gamma_{21}U_2+d_2V_2}{\Gamma_{21}^2+d_2^2},\qquad
\cos\gamma_{20}=\frac{\Gamma_{21}V_2-d_2U_2}{\Gamma_{21}^2+d_2^2}.
\end{align*}
These expressions follow directly from the steady mode-2 equations; for a consistent solution they must also satisfy $\sin^2\gamma_{20}+\cos^2\gamma_{20}=1$. On the symmetric branch, set $\Gamma_{21}=0$; the $\tilde{\eta}_{11}$ drive remains. Substitution into the mode-1 equations gives the explicit coupled response equation
\begin{equation}\label{eq:quadratic-response-explicit}
X_{12}^2+Y_{12}^2=\left(\frac{\mu}{2\hat{\omega}_1}\right)^2,
\end{equation}
where
\begin{align*}
X_{12}&=\frac{\hat{\zeta}_1}{2}a_{10}+\frac{\eta_{11}}{8}a_{10}^3+\frac{\Gamma_{12}}{4\hat{\omega}_1}a_{10}a_{20}\sin\gamma_{20},\\
Y_{12}&=\sigma_1a_{10}-\frac{3\kappa_{11}}{8\hat{\omega}_1}a_{10}^3-\frac{\kappa_{12}}{4\hat{\omega}_1}a_{10}a_{20}^2+\frac{\Gamma_{12}}{4\hat{\omega}_1}a_{10}a_{20}\cos\gamma_{20}.
\end{align*}
Together with the displayed sine/cosine formulas, \eqref{eq:quadratic-response-explicit} is a closed pair of amplitude equations for $(a_{10},a_{20})$; the external phase is recovered from $\cos\gamma_{10}=-2\hat{\omega}_1X_{12}/\mu$ and $\sin\gamma_{10}=-2\hat{\omega}_1Y_{12}/\mu$.

For the exact symmetric-branch reverse $2{:}1$ resonance, define
\[
P=\frac{q_{21}a_{10}}{4\hat{\omega}_2},\qquad
\Delta_{2,\mathrm{eff}}=\frac{\tilde{\sigma}_2+\sigma_1}{2}-\frac{3\kappa_{22}}{8\hat{\omega}_2}a_{20}^2-\frac{\kappa_{21}}{4\hat{\omega}_2}a_{10}^2.
\]
The steady mode-2 equations are $P\sin\tilde{\gamma}_{20}=\hat{\zeta}_2/2$ and $P\cos\tilde{\gamma}_{20}=-\Delta_{2,\mathrm{eff}}$. Eliminating $\tilde{\gamma}_{20}$ gives
\begin{equation}\label{eq:two-one-saturation}
\left(\frac{\hat{\zeta}_2}{2}\right)^2+\Delta_{2,\mathrm{eff}}^2=P^2.
\end{equation}
At onset, $a_{20}\to0$, $\gamma_{10}'=0$, and $\Delta_2=(\tilde{\sigma}_2+\sigma_1)/2-\kappa_{21}a_{10}^2/(4\hat{\omega}_2)$. Thus \eqref{eq:two-one-saturation} is the modal threshold condition. Since $\Delta_2$ contains $a_{10}^2$, the critical amplitude is generally implicit:
\[
\left(\frac{q_{21}a_{1,\mathrm{crit}}}{4\hat{\omega}_2}\right)^2=\left(\frac{\hat{\zeta}_2}{2}\right)^2+\left[\frac{\tilde{\sigma}_2+\sigma_1}{2}-\frac{\kappa_{21}a_{1,\mathrm{crit}}^2}{4\hat{\omega}_2}\right]^2.
\]
Equivalently, the exact threshold is the implicit amplitude equation
\[
a_{1,\mathrm{crit}}=\frac{2\hat{\omega}_2}{|q_{21}|}\sqrt{\hat{\zeta}_2^2+\left(\tilde{\sigma}_2+\sigma_1-\frac{\kappa_{21}}{2\hat{\omega}_2}a_{1,\mathrm{crit}}^2\right)^2}.
\]
Neglecting the amplitude-dependent cross-frequency shift gives $a_{1,\mathrm{crit}}=2\hat{\omega}_2|q_{21}|^{-1}\sqrt{\hat{\zeta}_2^2+(\tilde{\sigma}_2+\sigma_1)^2}$. The corresponding internal phase is $\sin\tilde{\gamma}_{20}=\hat{\zeta}_2/(2P)$ and $\cos\tilde{\gamma}_{20}=-\Delta_{2,\mathrm{eff}}/P$. Substituting these directly into the mode-1 balance, define
\begin{align*}
X_{21}&=\frac{\hat{\zeta}_1}{2}a_{10}+\frac{\eta_{11}}{8}a_{10}^3+\frac{a_{20}^2}{4\hat{\omega}_1P}\left(\frac{q_{12}\hat{\zeta}_2}{2}-\hat{\omega}_2\tilde{\eta}_{12}\Delta_{2,\mathrm{eff}}\right),\\
Y_{21}&=\sigma_1a_{10}-\frac{3\kappa_{11}}{8\hat{\omega}_1}a_{10}^3-\frac{\kappa_{12}}{4\hat{\omega}_1}a_{10}a_{20}^2-\frac{a_{20}^2}{4\hat{\omega}_1P}\left(q_{12}\Delta_{2,\mathrm{eff}}+\frac{\hat{\omega}_2\tilde{\eta}_{12}\hat{\zeta}_2}{2}\right),
\end{align*}
and require $X_{21}^2+Y_{21}^2=(\mu/(2\hat{\omega}_1))^2$. Together with \eqref{eq:two-one-saturation}, this gives the explicit coupled frequency-response equations for $(a_{10},a_{20})$.

For $1{:}3$ resonance, solve the mode-2 steady equations for the internal phase by defining
\begin{align*}
U_3&=\frac{4\hat{\omega}_2\hat{\zeta}_2a_{20}}{a_{10}^3},\\
V_3&=\frac{8\hat{\omega}_2a_{20}}{a_{10}^3}\left(\sigma_2-3\sigma_1+\frac{3\kappa_{22}}{8\hat{\omega}_2}a_{20}^2+\frac{\kappa_{21}}{4\hat{\omega}_2}a_{10}^2\right),\qquad d_3=\hat{\omega}_1\eta_{11},\\
\sin\gamma_{20}&=\frac{\Lambda_{21}U_3+d_3V_3}{\Lambda_{21}^2+d_3^2},\qquad
\cos\gamma_{20}=\frac{\Lambda_{21}V_3-d_3U_3}{\Lambda_{21}^2+d_3^2}.
\end{align*}
Again impose $\sin^2\gamma_{20}+\cos^2\gamma_{20}=1$ and substitute into
\begin{align*}
X_{13}&=\frac{\hat{\zeta}_1}{2}a_{10}+\frac{\eta_{11}}{8}a_{10}^3+\frac{\Lambda_{12}}{8\hat{\omega}_1}a_{10}^2a_{20}\sin\gamma_{20},\\
Y_{13}&=\sigma_1a_{10}-\frac{3\kappa_{11}}{8\hat{\omega}_1}a_{10}^3-\frac{\kappa_{12}}{4\hat{\omega}_1}a_{10}a_{20}^2+\frac{\Lambda_{12}}{8\hat{\omega}_1}a_{10}^2a_{20}\cos\gamma_{20}.
\end{align*}
The coupled frequency-response condition is $X_{13}^2+Y_{13}^2=(\mu/(2\hat{\omega}_1))^2$. For $1{:}1$, the corresponding explicit coupled equations are the four steady equations obtained from the trigonometric $m=1$ flow above; they are solved simultaneously for $(a_{10},a_{20},\gamma_{10},\gamma_{20})$ for each detuning.

To assess stability, linearize the real slow flow about a fixed point: $\delta\mathbf{u}'=\mathbf{J}_{\mathrm{SF}}\delta\mathbf{u}$, with $\delta\mathbf{u}=(\delta a_1,\delta\gamma_1,\delta a_2,\delta\gamma_2)^T$. At $a_{20}=0$, polar coordinates are singular. Use Cartesian envelope coordinates $A_k=\tfrac12(p_k-iq_k)e^{i\nu_kT_1}$ (equivalently, the real and imaginary parts of each complex envelope) to form the Jacobian. When the resonant forcing of mode 2 vanishes, it separates into the in-mode block $\mathbf{J}_{11}$ and transverse block $\mathbf{J}_{22}$. For the single-mode primary response, the in-mode block in polar coordinates is
\[
\mathbf{J}_{11}=\begin{pmatrix}
-\dfrac{\hat{\zeta}_1}{2}-\dfrac{3\eta_{11}}{8}a_{10}^2 & \dfrac{\mu}{2\hat{\omega}_1}\sin\gamma_{10}\\[4pt]
-\dfrac{3\kappa_{11}}{4\hat{\omega}_1}a_{10}-\dfrac{\mu}{2\hat{\omega}_1a_{10}^2}\sin\gamma_{10} & \dfrac{\mu}{2\hat{\omega}_1a_{10}}\cos\gamma_{10}
\end{pmatrix}.
\]
For the reverse $2{:}1$ resonance, the transverse mode-2 perturbation about the single-mode branch has parametric coefficient and half-detuning
\[
P=\frac{q_{21}a_{10}}{4\hat{\omega}_2},\qquad
\Delta_2=\frac{\tilde{\sigma}_2+\sigma_1}{2}-\frac{\kappa_{21}}{4\hat{\omega}_2}a_{10}^2,
\]
where $\gamma_{10}'=0$ on that branch. In Cartesian coordinates obtained by resolving the half-relative phase, the transverse Jacobian is
\[
\mathbf{J}_{22}=\begin{pmatrix}
-\hat{\zeta}_2/2&-\Delta_2+P\\
\Delta_2+P&-\hat{\zeta}_2/2
\end{pmatrix},\qquad
\det(\mathbf{J}_{22})=\frac{\hat{\zeta}_2^2}{4}+\Delta_2^2-P^2.
\]
Thus the zero-eigenvalue threshold is
\[
\left(\frac{q_{21}a_{10}}{4\hat{\omega}_2}\right)^2
=\left(\frac{\hat{\zeta}_2}{2}\right)^2
+\left(\frac{\tilde{\sigma}_2+\sigma_1}{2}-\frac{\kappa_{21}}{4\hat{\omega}_2}a_{10}^2\right)^2.
\]

On a coupled branch, let $c_1=-\operatorname{tr}\mathbf{J}_{\mathrm{SF}}$, $c_2$ be the sum of its six principal $2\times2$ minors, $c_3$ the negative sum of its four principal $3\times3$ minors, and $c_4=\det\mathbf{J}_{\mathrm{SF}}$. Its characteristic polynomial is
\begin{equation}\label{eq:slowflow-characteristic}
P(\nu)=\nu^4+c_1\nu^3+c_2\nu^2+c_3\nu+c_4.
\end{equation}
The Routh--Hurwitz conditions for all four slow-flow eigenvalues to have negative real parts are
\[
c_1>0,\quad c_2>0,\quad c_3>0,\quad c_4>0,\qquad
\Delta_3:=c_1c_2c_3-c_1^2c_4-c_3^2>0.
\]
The condition $c_4=0$ is a zero-eigenvalue (generically saddle-node) boundary, while $\Delta_3=0$ with the other inequalities strict is the Hopf boundary. A slow-flow Hopf bifurcation produces modulation of the response envelopes and hence quasiperiodic or otherwise nonstationary motion in the original oscillator.

\subsubsection{Frequency Response}
\begin{assignmentinstruction}
*Plot the amplitudes as a function of the external forcing frequency (external excitation near primary resonance OR parametric excitation near the principle parametric resonance).
\end{assignmentinstruction}
For fixed physical parameters, sweep $\sigma_1$ (equivalently $\Omega=\hat{\omega}_1+\varepsilon\sigma_1$), solve the steady slow-flow equations at each point, and continue both the single-mode response \eqref{eq:single-mode-response} and any coupled branches. Classify each point using the eigenvalues of $\mathbf{J}_{\mathrm{SF}}$; stable and unstable portions are conventionally distinguished by solid and dashed curves, respectively. The coupled amplitudes depend on both detunings through the phase-locked relations above. No response plot is inserted here.
% TODO: Insert Frequency Response Figure here

\section{Closing Remarks}
\begin{assignmentinstruction}
Summarize your results and conclusions from: (i) equilibrium stability and the conditions for internal resonance (ii) frequency response of the system due to external OR parametric excitation. Comment on possible existence of quasiperiodic or non-stationary response (e.g. there exists a Hopf bifurcation in the slowly varying evolution equations).
\end{assignmentinstruction}

\section{Optional Bonus: Numerical Validation}
\begin{assignmentinstruction}
Verify your findings numerically using a fourth-order Runge-Kutta ODE solver [e.g. Matlab - ODE45 with sufficiently small tolerances $\sim O(10^{-8})$]: (i) solution of the slowly varying evolution equations vs (ii) solution of the original system.

Comment on accuracy of the (approximate) analytical maximal amplitudes in the frequency response function deduced form the slowly varying evolution equations in comparison to selected maximal amplitudes of the periodic response in the original sleeping gyro equations of motion.
\end{assignmentinstruction}

\begin{thebibliography}{1}
\bibitem{alcorta2021}
R. Alcorta, A. Geminterna, R. H. Blick, and O. Gottlieb,
``Nonlinear interactions in a hybrid multi-field nanoresonator subject to combined external and parametric excitation,'' 2021.
\end{thebibliography}

\end{document}
